\section{Conclusions}
\label{sec:conclusions}

FilaSeg reconstructs overlapping filament instances from degraded centreline
masks through deterministic Hough vectorization, orientation-aware cleaning and
geometric reconnection. On \LockedSceneCount{} predefined procedural scenes, its
Stage-3 centreline output achieved common-fragment \Fone{} of
\LockedFilaSegCommonFOneMean{} (density-stratified 95\% CI
[\LockedFilaSegCommonFOneCILow{}, \LockedFilaSegCommonFOneCIHigh{}]). The paired
difference from the fixed minimum-turn comparator was
\FilaSegMinusBaselineSkeletonCommonFOneMean{} (95\% CI
[\FilaSegMinusBaselineSkeletonCommonFOneCILow{},
\FilaSegMinusBaselineSkeletonCommonFOneCIHigh{}]).

The evaluator audit changed the scientific interpretation without changing any
frozen method output. Equivalent Stage-3 centrelines are the defensible primary
grouping representation; Stage-4 centreline changes and width-rendered masks are
secondary diagnostics. Automatically selected counterexamples show that the
aggregate advantage is not uniform across scenes or densities. The result is
therefore evidence for this fixed comparison under the tested generator, not a
claim of superiority over filament tracing methods generally.

The single real crop remains a development illustration. Independent real
annotations from multiple parent images, complete acquisition and annotation
metadata, separate validation of apparent width and other exported
measurements, and a citable clean-checkout release are required before broader
SEM or quantitative microstructure claims are justified.
